\chapter{Coupon Collector's Test}

\section{The Basic Story}
The coupon collector's problem\cite{wiki:coupon-collector} is a classical question in probability.  Imagine
a ``collect all coupons and win'' promotion: each cereal box contains one
random coupon, chosen uniformly from \(n\) different designs.  You want to know:
\begin{quote}
  \emph{How many boxes will I need to buy, on average, before I have seen
  every coupon at least once?}
\end{quote}

Equivalently: draw tickets with replacement from a set of \(n\) labels
\(\{1,2,\dots,n\}\).  Let \(T\) be the number of draws needed until you have
seen all \(n\) labels.  The coupon collector's problem studies the distribution
of \(T\), in particular its \emph{expectation} and \emph{variance}.\index{Coupon collector's problem}

\section{Expectation: How Long Do We Wait on Average?}
We build up the collection step by step.  Suppose you have already collected
\((i-1)\) distinct coupons.  The chance that the next box contains a \emph{new}
coupon (one of the remaining \(n-(i-1)\)) is
\[
  p_i = \frac{n-(i-1)}{n} = \frac{n-i+1}{n}.
\]
The number of additional boxes \(t_i\) you need to buy in order to get the
\(i\)-th new coupon has a geometric distribution with mean
\[
  \mathbb{E}[t_i] = \frac{1}{p_i} = \frac{n}{n-i+1}.
\]

The total time to collect all \(n\) coupons is
\[
  T = t_1 + t_2 + \cdots + t_n,
\]
so by linearity of expectation,
\[
  \mathbb{E}[T]
  = \mathbb{E}[t_1] + \cdots + \mathbb{E}[t_n]
  = \frac{n}{n} + \frac{n}{n-1} + \cdots + \frac{n}{1}
  = n \left( 1 + \frac{1}{2} + \cdots + \frac{1}{n} \right)
  = n\,H_n,
\]
where \(H_n\) is the \(n\)-th harmonic number.  Using the asymptotic expansion
for \(H_n\),
\[
  H_n = \log n + \gamma + \frac{1}{2n} + O\!\left(\frac{1}{n^2}\right),
\]
we obtain the intuitive rule of thumb
\[
  \mathbb{E}[T] \approx n \log n + \gamma n + \frac{1}{2},
\]
where \(\gamma \approx 0.5772\) is the Euler--Mascheroni constant.\index{Harmonic numbers}

The important qualitative message is that the expected time grows like
\(n \log n\): much longer than the simple \(n\) you might guess from ``one of
each''.  The last few missing coupons take disproportionately long to find.

\section{Variance: How Much Can the Waiting Time Fluctuate?}
Because the \(t_i\) are independent, we can add their variances:
\[
  \mathrm{Var}(T)
  = \mathrm{Var}(t_1) + \cdots + \mathrm{Var}(t_n).
\]
For a geometric random variable with success probability \(p_i\),
\[
  \mathrm{Var}(t_i) = \frac{1-p_i}{p_i^2}.
\]
Substituting \(p_i = (n-i+1)/n\) gives
\[
  \mathrm{Var}(T)
  = \frac{1-p_1}{p_1^2} + \cdots + \frac{1-p_n}{p_n^2}
  = n^2\left(\frac{1}{1^2} + \frac{1}{2^2} + \cdots + \frac{1}{n^2}\right)
    - n\left(1 + \frac{1}{2} + \cdots + \frac{1}{n}\right).
\]
Using the fact that
\[
  \sum_{k=1}^{\infty}\frac{1}{k^2} = \frac{\pi^2}{6},
\]
one obtains the bound
\[
  \mathrm{Var}(T) < \frac{\pi^2}{6}\,n^2.
\]

This tells us that the spread of \(T\) around its mean is on the order of
\(n\): even though the expected time is about \(n\log n\), for a fixed \(n\) we
should not be surprised if the actual time fluctuates by a multiple of \(n\)
either way.

\section{Exact Distribution: All Coupons Seen After \texorpdfstring{$m$}{m} Draws}
So far we have focused on the mean and variance of \(T\), the number of draws
needed to see all \(n\) coupons.  We can also ask for the \emph{exact
probability} that after exactly \(m\) draws, we have already seen all \(n\)
coupons at least once.

Think of a specific sequence of \(m\) draws as a list
\[
  (c_1,c_2,\dots,c_m),
\]
where each \(c_j\) is one of the coupon types \(1,2,\dots,n\).  There are
exactly \(n^m\) possible such lists, since each position can contain any of the
\(n\) coupons.

We want to count how many of these lists contain \emph{every} coupon type at
least once.  A direct count is hard, but an inclusion--exclusion argument
gives a clean formula (see, for example,\cite{693222-combinatorics-coupon}).

\begin{itemize}
  \item Start with all \(n^m\) lists.
  \item Subtract the lists that miss at least one coupon type.  For each choice
        of a particular coupon to be missing, there are \((n-1)^m\) lists that
        never use it, so this would suggest subtracting
        \(\binom{n}{1}(n-1)^m\).
  \item But now we have subtracted too much: lists that miss at least two
        coupon types have been subtracted once for each missing coupon.  So we
        add back, for each pair of coupons, the \((n-2)^m\) lists that use only
        the remaining \(n-2\) coupons, giving a term
        \(\binom{n}{2}(n-2)^m\).
  \item Continuing this inclusion--exclusion process, we alternate minus and
        plus signs, moving down to lists that use only \(n-k\) coupon types.
\end{itemize}

The final count of length-\(m\) draw sequences that contain all \(n\) coupon
types at least once is
\[
  \sum_{k=0}^{n} (-1)^k \binom{n}{k}\,(n-k)^m.
\]
Since all \(n^m\) sequences are equally likely, the probability that \(m\)
draws already contain all \(n\) coupons is
\[
  P(\text{all $n$ coupons seen in $m$ draws})
  = \frac{1}{n^m}
    \sum_{k=0}^{n} (-1)^k \binom{n}{k}\,(n-k)^m.
\]
This gives the exact distribution of the event ``after \(m\) draws we have
collected the full set'', for any fixed \(m\) and \(n\).

\section{Stirling Numbers of the Second Kind}
Stirling numbers of the second kind\cite{wiki:stirling-second-kind} arise
throughout combinatorics.  They count the number of ways to partition a set of
\(n\) labelled objects into \(k\) non-empty, unlabeled subsets.  The standard
notations are
\[
  S(n,k) \quad \text{or} \quad \left\{\begin{matrix}n\\k\end{matrix}\right\}.
\]
For example, there are
\(\left\{\begin{smallmatrix}4\\2\end{smallmatrix}\right\}=7\) ways to split a
4-element set into 2 non-empty groups: \\
(1), (2,3,4)\\
(2), (1,3,4)\\
(3), (1,2,4)\\
(4), (1,2,3)\\
(1,2), (3,4)\\
(1,3), (2,4)\\
(1,4), (2,3).

\section{Connect Coupon Collector's Problem to Stirling Numbers of Second Kind}
We shall prove that number of sequences that \(n\) draws contains \(k\) different values 
equals to 
the number of ways to partition a set of
\(n\) labelled objects into \(k\) non-empty, \emph{labeled} subsets:

A draw $(a_1, a_2, \dots, a_m)$ with each $a_i \in \{1,\dots,n\}$ induces a partition in which
the subset label is the value of $a_i$ and the object label is the sequence index $i$.

For example, the draw $(1,3,1,2)$ induces the partition
\[
  1:\{1,3\},\quad 2:\{4\},\quad 3:\{2\}.
\]

Therefore the number of different draws that have \(n\) coupons and \(k\) different values equals to 
the number of ways to partition a set of
\(n\) labelled objects into \(k\) non-empty, \emph{labeled} subsets.

Stirling numbers of the second kind count the number of ways to partition a set of
\(n\) labelled objects into \(k\) non-empty, \emph{unlabeled} subsets. 
Therefore number of different draws that have n coupons and\(k\) different values equals to 
\(S(n,k) \cdot k!\).

\section{Coupon Collector's Test for PowerSpin}
Our basic coupon collector model assumes that all coupon types are equally
likely.  In PowerSpin, however, a symbol is three times as likely as any single
number, so we must first re-define what we mean by a ``coupon type''.

Group the wheel positions into nine types as follows: for \(k=1,\dots,8\), the
numbers \(3k-2,3k-1,3k\) form coupon type \(k\), and all symbol positions form
coupon type \(9\).  Under the official rules, each of these nine types then has
the same probability.

With this redefinition, the coupon collector formulas above (expressed via
Stirling numbers of the second kind) give the theoretical distribution of the
number of spins needed to see all nine types.  From the PowerSpin data we can
estimate the empirical distribution, and a chi-squared test can be used to
check whether the observed behaviour is consistent with the theoretical model.
